\documentclass[11pt]{article}

% ------------------------------------------------------------
% Packages
% ------------------------------------------------------------
\usepackage[margin=1in]{geometry}
\usepackage{amsmath, amssymb, amsthm, mathtools}
\usepackage{graphicx}
\usepackage{booktabs}
\usepackage{hyperref}
\usepackage{enumitem}
\usepackage{physics}
\usepackage{authblk}
\usepackage{setspace}
\usepackage{microtype}

\onehalfspacing

% ------------------------------------------------------------
% Theorem environments
% ------------------------------------------------------------
\newtheorem{theorem}{Theorem}
\newtheorem{definition}{Definition}
\newtheorem{proposition}{Proposition}
\newtheorem{corollary}{Corollary}
\newtheorem{remark}{Remark}

% ------------------------------------------------------------
% Custom commands
% ------------------------------------------------------------
\newcommand{\CVT}{\mathrm{CVT}}
\newcommand{\V}{\mathcal{V}}
\newcommand{\Hhost}{\mathcal{H}}
\newcommand{\Rhost}{R_{\mathrm{host}}}
\newcommand{\RP}{R_P}
\newcommand{\Phii}{\Phi_i}
\newcommand{\Jraw}{J_{\mathrm{raw}}}
\newcommand{\Jeff}{J_{\mathrm{eff}}}
\newcommand{\chialpha}{\chi_{\alpha}}

% ------------------------------------------------------------
% Title
% ------------------------------------------------------------
\title{\textbf{Coherence Viability Theory:}\\
Bounded Exchange, Receptive Basins, and the Conditions for Hosted Transformation}

\author{Daniel Varela Arevalo}
\date{2026}

\begin{document}

\maketitle

\begin{abstract}
We propose Coherence Viability Theory (CVT) as a testable framework for boundary-mediated transformations in which the continued viability of an identifiable receiving or hosting system is part of the success criterion. CVT hypothesizes a candidate non-compensatory structure comprising bounded coupling, preserved identity-defining organization, pre-existing capture readiness, and available restorative reserve. These elements are not individually new: the framework draws from viability theory, autopoiesis, passivity-based control, resilience theory, and basin dynamics. Its proposed contribution is their conjunction as conditions of hosted transformation, together with a distinction between exchange delivered to a system and exchange retained as productive uptake. The non-compensability claim is conceptually prior to any particular scalar aggregator; the product-gate form developed here is one candidate representation rather than an established universal equation. We develop a minimal gated architecture and a membrane formulation, then use artificial fusion as an illustrative toy-model instantiation of the difference between initiating a transition and hosting it. CVT does not yet establish that the proposed conditions are universally necessary or sufficient. It instead defines a bounded research program in which the conditions, observables, aggregation rules, and counterexamples can be compared empirically or in reproducible simulation.
\end{abstract}

\section{Introduction}

Many systems depend on exchange across a boundary. They must receive energy, matter, information, influence, or control while retaining the organization that allows them to continue functioning as hosts. Neither closure nor openness alone solves this problem. Insufficient exchange can starve a system; excessive or poorly received exchange can overload it, erase relevant distinctions, exhaust recovery capacity, or drive it outside an admissible regime.

This paper studies a bounded class of such processes: \emph{boundary-mediated transformations in which continued host viability is part of what counts as success}. We call a transformation \emph{hosted} when a declared host undergoes or supports a meaningful transition while remaining viable over a stated observation horizon. The host, boundary, transformation variable, success criterion, timescale, and level of analysis must therefore be specified before CVT is applied.

CVT proposes that hosted transformation may depend on four jointly necessary conditions:
\begin{enumerate}[label=(\roman*)]
    \item \textbf{bounded coupling}: exchange is present but constrained relative to the host's changing handling range;
    \item \textbf{preserved distinction}: the host retains the identity-defining organization or independent constraints relevant to the declared success criterion;
    \item \textbf{capture readiness}: before uptake is observed, the host occupies a state from which delivered input can enter and remain within an admissible basin;
    \item \textbf{available restorative reserve}: sufficient internal or reliably accessible recovery capacity remains to absorb strain, repair damage, and preserve control authority.
\end{enumerate}

In compact form, the candidate structure is
\begin{equation}
    \mathrm{CVT}_{\mathrm{candidate}}
    =
    B\cap Q\cap C\cap S.
\end{equation}
This expression is a hypothesis about essential conditions, not a proof that the conditions are universal or sufficient. Their satisfaction defines a candidate viability region. Hidden state variables, delayed damage, environmental change, path dependence, adversarial interference, and measurement error may still produce failure.

A central distinction is between what is delivered and what is productively retained. Let \(J_{\mathrm{raw}}\) denote boundary flux delivered to the host and \(J_{\mathrm{eff}}\) denote retained productive uptake over a declared horizon. Realized uptake efficiency may be observed as
\begin{equation}
    \rho_{\mathrm{obs}}
    =
    \frac{J_{\mathrm{eff}}}{J_{\mathrm{raw}}},
\end{equation}
when \(J_{\mathrm{raw}}>0\). This ratio is an outcome, not an explanation. CVT must predict it from independently measured host conditions rather than define receptivity through the same gates later used to explain it. The testable claim is that increased delivered exchange can, in some regimes, reduce productive uptake when capture readiness or restorative reserve degrades sufficiently.

The individual ingredients of this framework are inherited from mature traditions. Viability theory studies admissible trajectories under state constraints. Autopoiesis studies self-producing organization and structural coupling. Passivity and port-Hamiltonian control formalize constrained energetic exchange. Resilience theory studies persistence and recovery under disturbance. Basin dynamics studies nonlinear capture and escape. CVT does not claim these concepts as new. Its proposed contribution is narrower: it asks whether a class of boundary-mediated failures is better understood when bounded coupling, identity preservation, capture readiness, and restorative reserve are treated as non-compensatory rather than substitutable.

Non-compensability is conceptually prior to the product equation used later in the manuscript. A product of normalized gates is one smooth representation because a near-zero essential gate sharply lowers the aggregate. It is not yet established as the uniquely correct form. Minimum-gate, weighted geometric-mean, normalized soft-minimum, and simultaneous constraint-intersection formulations remain competing representations to be compared against the same data or simulations.

The membrane formulation provides the paper's clearest structural illustration. A boundary can permit flux while the receiving domain remains unable to retain that flux productively. Permeability is therefore not identical to capture readiness or effective uptake. The fusion section extends this logic as an illustrative toy model: ignition access is distinguished from hosted burn, but the numerical thresholds and controller comparisons are not presented as calibrated reactor-physics results or as validation of CVT.

The paper proceeds as follows. Section~\ref{sec:lineage} situates CVT within its principal intellectual lineages. Section~\ref{sec:invariant} states the candidate four-condition structure and the non-compensability proposal. Section~\ref{sec:formal} develops an initial gated formal architecture whose variables and aggregation rules remain subject to operational testing. Section~\ref{sec:membrane} gives the membrane formulation. Section~\ref{sec:fusion} presents the fusion material as an illustrative instantiation of hosted transformation. The present paper establishes neither universal necessity nor sufficiency; it defines a research framework whose strongest claims must survive counterexamples, independent measurement, competing models, and delayed-failure tests.

\section{Intellectual Lineage}
\label{sec:lineage}

CVT is assembled from established literatures on viability, self-maintaining organization, open-system control, resilience, and nonlinear basin structure. The present claim is not that these traditions unknowingly contain one already-demonstrated CVT invariant. It is narrower: their concepts motivate a proposed decomposition of hosted transformation into bounded coupling, preserved identity-defining organization, pre-existing capture readiness, and available restorative reserve. Whether that decomposition adds explanatory or predictive value remains a question for comparison, intervention, and counterexample testing.

\subsection{Viability Theory}

Viability theory studies evolutions constrained to remain within admissible sets and characterizes the initial states from which at least one viable trajectory exists, including through viability kernels and regulation under constraints \cite{aubin2011}. CVT inherits this concern with admissible evolution. Its proposed emphasis is boundary-mediated forcing: the host must not only remain inside constraints, but do so while receiving inputs whose rate, timing, and form may alter the host's capacity to remain viable.

This is not yet a mathematical extension of viability theory. It is a proposed decomposition of one class of viability problems. CVT earns distinct value only if the explicit separation of boundary regulation, capture readiness, identity preservation, and restorative reserve improves analysis beyond an ordinary state-constraint formulation.

\subsection{Autopoiesis and Structural Coupling}

Maturana and Varela describe living systems through autopoietic organization: a network of processes that continuously produces and maintains the components and boundary constituting the system as a unity \cite{maturana1980}. CVT inherits the importance of organization and boundary-mediated relation, but it does not identify every CVT host as autopoietic and does not claim organizational closure or structural coupling as new.

Preserved distinction in CVT therefore means continuity of the identity-defining organization relevant to the declared host and success criterion. It does not mean static form, impermeability, or isolation. A host may change profoundly while retaining the organization or constraints that make it the same host for the purpose of the analysis.

\subsection{Passivity and Port-Hamiltonian Systems}

Port-Hamiltonian theory models open physical systems through storage, interconnection, dissipation, and ports that expose exchanges with an environment \cite{vanderschaft2014}. It provides a substantially more rigorous account of energetic exchange than CVT presently offers. CVT does not modify that formalism in this manuscript.

The proposed CVT question is whether exchange capacity and productive retention must be represented separately in some hosted transformations, and whether identity, capture, or repair constraints matter in addition to energetic balance. In physical applications, those claims must be translated into domain-native variables and tested against port-Hamiltonian, passivity, and constraint-based baselines rather than asserted by analogy.

\subsection{Resilience and Recovery}

Holling distinguished ecological resilience from local stability by emphasizing persistence and the capacity of systems to absorb change without losing their organizing relations \cite{holling1973}. CVT inherits this concern with persistence under disturbance. Its term \emph{available restorative reserve} is intended to be narrower than resilience as a whole: it denotes currently accessible recovery headroom, not every property that contributes to long-run persistence.

Restorative reserve may coincide with an established resilience measure in some domains and add nothing in others. It should be retained as a separate CVT variable only where it can be independently operationalized and improves prediction of recovery, delayed damage, or repeated-load failure.

\subsection{Basin Dynamics and Capture Readiness}

Basin-based approaches characterize nonlocal stability by considering the set or measure of initial conditions that return to or remain associated with a desired attractor after finite perturbation. Basin stability, for example, complements local linear stability with a nonlocal measure related to basin volume \cite{menck2013}.

CVT uses this lineage to define \emph{capture readiness}: before uptake is observed, the host occupies a state from which a specified input can enter and remain within an admissible basin over a declared horizon. Capture readiness is not the same as the later outcome of productive uptake. A boundary may deliver input into a system that has no suitable capture path, while a receptive state may remain inaccessible because the boundary does not deliver the required input.

\subsection{Status of the Synthesis}

These traditions overlap, but conceptual overlap does not establish a shared mechanism or a universal invariant. CVT's present contribution is a testable synthesis hypothesis:

\begin{quote}
For a scoped class of boundary-mediated hosted transformations, failure may be better explained when bounded coupling, identity preservation, capture readiness, and restorative reserve are treated as jointly required and potentially non-compensatory.
\end{quote}

This claim is supported only when the CVT decomposition can be operationalized independently, compared with simpler or domain-native models, and exposed to cases in which one proposed condition is absent without loss of hosted viability.

\section{Candidate CVT Structure and Non-Compensability}
\label{sec:invariant}

\begin{definition}[Hosted coherence viability]
Fix a declared host, boundary, transformation, success criterion, observation horizon, and level of analysis. The host is coherence-viable for that transformation when it can undergo or support the transition while remaining inside the declared viable set and retaining the identity-defining organization required by the success criterion.
\end{definition}

Let \(\mathcal K_B(t)\), \(\mathcal K_Q(t)\), \(\mathcal K_C(t)\), and \(\mathcal K_S(t)\) denote the state sets satisfying four proposed conditions at time \(t\):

\begin{enumerate}[label=(\roman*)]
    \item \textbf{bounded coupling} \(B\): the exchange schedule remains within a domain-specific admissible range relative to the host's changing handling capacity;
    \item \textbf{preserved identity-defining organization} \(Q\): the organization, constraints, or independent control structure relevant to the declared host persists through the transformation;
    \item \textbf{capture readiness} \(C\): before realized uptake is measured, the host state permits the specified input to enter and remain within an admissible basin over the stated horizon;
    \item \textbf{available restorative reserve} \(S\): sufficient internal or reliably accessible recovery capacity remains to absorb strain, repair damage, or preserve control authority.
\end{enumerate}

The candidate CVT viability region is
\begin{equation}
    \mathcal K_{\CVT}(t)
    =
    \mathcal K_B(t)
    \cap
    \mathcal K_Q(t)
    \cap
    \mathcal K_C(t)
    \cap
    \mathcal K_S(t).
    \label{eq:candidate-cvt-region}
\end{equation}
The shorthand \(B\cap Q\cap C\cap S\) refers to simultaneous satisfaction of these conditions; it should not be read as though four scalar variables were themselves sets.

Equation~\eqref{eq:candidate-cvt-region} states a candidate structure, not a demonstrated universal law. CVT presently hypothesizes that the four conditions may be necessary in some scoped domains. It does not establish that every condition is universally necessary, and their simultaneous satisfaction is not presently sufficient for success. Delayed damage, hidden variables, regime change, path dependence, adversarial interference, and measurement error may still produce failure.

\subsection{Conceptual Non-Compensability}

A proposed condition is non-compensatory in a specified domain when crossing its failure threshold prevents the declared viable outcome despite admissible improvement in the remaining conditions, unless the failed condition itself is restored or the scope of the host is changed. In gate notation, a strong form would be
\begin{equation}
    G_k < \Theta_k
    \quad\Longrightarrow\quad
    y \notin \mathcal Y_{\mathrm{viable}}
    \qquad
    \text{for all admissible values of }\{G_i\}_{i\neq k},
    \label{eq:strong-noncompensability}
\end{equation}
under stated domain assumptions.

Equation~\eqref{eq:strong-noncompensability} is an empirical or model-class hypothesis. A valid case in which severe failure of one condition is reliably offset by another would constrain strict non-compensability, reveal an omitted restoration pathway, or show that the failed condition was not essential in that domain.

\subsection{Product Gate as a Candidate Representation}

One smooth representation of essential-gate fragility is
\begin{equation}
    \V_{\Pi}
    =
    \prod_{i=1}^{n}G_i,
    \qquad G_i\in[0,1].
    \label{eq:product-candidate}
\end{equation}

\begin{proposition}[Zero-gate property of the product representation]
If viability is represented by Equation~\eqref{eq:product-candidate}, then \(G_k\to 0\) implies \(\V_{\Pi}\to 0\), regardless of the remaining gate values.
\end{proposition}

\begin{proof}
Because each \(G_i\in[0,1]\),
\[
    \V_{\Pi}
    =
    G_k\prod_{i\neq k}G_i.
\]
The remaining product is bounded above by one, so \(G_k\to0\) implies \(\V_{\Pi}\to0\).
\end{proof}

This proposition proves a property of the selected aggregator; it does not prove that real systems aggregate viability multiplicatively. The product must therefore compete with alternatives such as
\begin{equation}
    \V_{\min}=\min_i G_i,
    \qquad
    \V_G=\prod_i G_i^{w_i},\quad \sum_i w_i=1,
\end{equation}
and with direct simultaneous constraint satisfaction \(G_i\geq\Theta_i\) for every essential condition. Model comparison should use independently defined viability outcomes, common data or simulations, and out-of-sample evaluation.

\subsection{Capture Readiness and Realized Uptake}

Capture readiness \(C\) is a pre-existing state condition. Realized uptake efficiency is an observed outcome:
\begin{equation}
    \rho_{\mathrm{obs}}
    =
    \frac{J_{\mathrm{eff}}}{J_{\mathrm{raw}}},
    \qquad J_{\mathrm{raw}}>0.
\end{equation}
A model that defines \(C\) by \(\rho_{\mathrm{obs}}\) cannot then use \(C\) to explain the same observed ratio. CVT must operationalize capture readiness from state, basin, or response-capacity information available before or independently of the uptake outcome.

The immediate research burden is therefore not to defend one equation. It is to determine, domain by domain, whether the four proposed conditions are identifiable, whether any are genuinely essential, and which mathematical representation best predicts hosted viability without circular measurement.

% CVT Foundations v0.2 — proposed Formal Architecture replacement
% Review artifact. This section replaces only the material between
% \section{Formal Architecture} and \section{Membrane Formulation}.

\section{Formal Architecture: Observables, Gates, and Candidate Models}
\label{sec:formal}

The formal architecture must separate what is measured from how it is normalized and from how normalized quantities are combined. CVT therefore distinguishes three layers:
\begin{enumerate}[label=(\roman*)]
    \item \textbf{observable layer}: domain-native measurements with units, time windows, and uncertainty;
    \item \textbf{gate layer}: documented transformations of those measurements into normalized scores;
    \item \textbf{model layer}: candidate aggregation and dynamical rules whose performance can be compared against an independently defined viability outcome.
\end{enumerate}
A normalized gate is not itself an observation, and no scalar aggregator defines success merely by being chosen.

\subsection{Observable and Gate Layers}

For each proposed condition \(i\in\mathcal I=\{B,Q,C,S\}\), let \(z_i(t)\) denote a raw domain-specific measurement or vector of measurements. A normalization map
\begin{equation}
    G_i(t)
    =
    g_i\!\left(z_i(t);\vartheta_i\right)
    \in[0,1]
    \label{eq:gate-normalization}
\end{equation}
converts the raw measurement into a gate score. The parameters \(\vartheta_i\), normalization range, uncertainty interval, and threshold rationale must be reported for each application.

The maps need not be monotone. In particular, bounded coupling generally requires a two-sided map: both insufficient exchange and excessive exchange may receive low scores, while an admissible operating band receives a high score. Preserved organization \(Q\), capture readiness \(C\), and available restorative reserve \(S\) likewise require domain-native observables rather than intuitive numerical assignments.

Let \(y(t)\) denote a domain-native viability outcome and let \(\mathcal Y_{\mathrm{viable}}\) be declared before the gates are fitted. The binary event
\begin{equation}
    V_{\mathrm{set}}(t)
    =
    \mathbf{1}\!\left\{y(t)\in\mathcal Y_{\mathrm{viable}}\right\}
    \label{eq:independent-viability-event}
\end{equation}
is an independently defined outcome, not a CVT aggregate. Continuous domain-native outcomes may be used instead when binary classification would discard relevant information.

\subsection{Candidate Aggregation Family}

The intersection in Equation~\eqref{eq:candidate-cvt-region} specifies simultaneous constraint satisfaction. When a scalar score is useful, several non-compensatory or weakly compensatory representations should be compared:
\begin{align}
    \V_{\Pi}(t)
    &=
    \prod_{i\in\mathcal I}G_i(t),
    \\
    \V_{\min}(t)
    &=
    \min_{i\in\mathcal I}G_i(t),
    \\
    \V_G(t)
    &=
    \prod_{i\in\mathcal I}G_i(t)^{w_i},
    \qquad
    w_i\geq0,
    \quad
    \sum_iw_i=1,
    \\
    \V_{\tau}(t)
    &=
    -\tau\log\!\left(
        \frac{1}{|\mathcal I|}
        \sum_{i\in\mathcal I}e^{-G_i(t)/\tau}
    \right),
    \qquad \tau>0.
    \label{eq:aggregation-family}
\end{align}
The last expression is a normalized soft minimum. Direct simultaneous constraints \(G_i\geq\Theta_i\) for every essential condition remain a further alternative. Product, minimum, weighted geometric mean, soft minimum, and constraint intersection encode different assumptions and must be evaluated against the same independently defined outcome using common data or simulations and out-of-sample comparison.

\subsection{Threshold Violations}

Given a gate threshold \(\Theta_i\), define the one-sided violation
\begin{equation}
    \Phi_i(t)
    =
    [\Theta_i-G_i(t)]_+,
    \qquad
    [x]_+=\max\{x,0\}.
    \label{eq:violation-variable}
\end{equation}
A two-sided raw constraint should be encoded in the normalization map \(g_i\), so that \(\Phi_i\) retains one interpretation: distance below the normalized viability threshold. Thresholds should be justified independently, preregistered when possible, or estimated out of sample. The uncertainty in \(\Phi_i\) inherits uncertainty in both \(z_i\) and \(\Theta_i\).

\subsection{Candidate Gate Dynamics}

A phenomenological recovery model may be written as
\begin{equation}
    \dot G_i
    =
    a_i(H,S)\Phi_i
    -
    \ell_i(t)
    -
    \sum_{j\neq i}\kappa_{ij}\Phi_j,
    \label{eq:gate-dynamics}
\end{equation}
with projected or saturating dynamics used to preserve \(G_i\in[0,1]\). Here \(\ell_i(t)\) is a normalized gate-degradation rate with units of inverse time, not an unscaled physical load. The coefficients \(\kappa_{ij}\) represent lagged cross-gate degradation and cannot be inferred from simultaneous decline alone, because common-cause forcing is an alternative explanation.

One candidate recovery coefficient is
\begin{equation}
    a_i(H,S)
    =
    \alpha_i(1+\eta_i H)q_i(S),
    \label{eq:recovery-coefficient}
\end{equation}
where \(q_i\) is nondecreasing and may satisfy \(q_i(0)=0\) when recovery requires available reserve. Equation~\eqref{eq:gate-dynamics} is not a universal law. It is a model class whose terms must be mapped to domain-native mechanisms and compared with simpler baselines.

\subsection{Recovery History and Available Reserve}

Recovery history or maturity and present restorative reserve are distinct states. Let
\[
    H(t)\in[0,1]
\]
denote learned recovery history, conditioning, or maturity, and let
\[
    S(t)\in[0,S_{\max}]
\]
denote currently available restorative reserve. The normalized reserve gate is \(G_S=g_S(S)\).

A candidate learning law is
\begin{equation}
    \dot H
    =
    \epsilon U_{\mathrm{rec}}(t)(1-H)
    -
    \zeta H
    -
    \xi_D D H,
    \label{eq:history-dynamics}
\end{equation}
where \(U_{\mathrm{rec}}(t)\) is an independently coded measure of successful viable operation or recovery. A candidate reserve law is
\begin{equation}
    \dot S
    =
    r_S(H)(S_{\max}-S)
    -
    c_S(t)
    -
    \sum_i\nu_i\Phi_i,
    \label{eq:reserve-dynamics}
\end{equation}
with projection onto \([0,S_{\max}]\). The term \(c_S(t)\) is measured reserve consumption; \(r_S(H)\) is a replenishment rate that may improve with history. These equations permit a mature system to be temporarily exhausted and an inexperienced system to possess substantial unused reserve. Neither possibility was representable when memory, maturity, and reserve shared one variable.

\subsection{Raw Exchange, Productive Uptake, and Prediction}

Let \(J_{\mathrm{raw}}(t)\) be measured delivered exchange. In a transport model one may have
\begin{equation}
    J_{\mathrm{raw}}
    =
    \Pi\Delta,
    \label{eq:raw-exchange}
\end{equation}
where \(\Pi\) is permeability or exchange capacity and \(\Delta\) is a driving gradient. This relation is domain-specific rather than the definition of raw exchange in every application.

Let \(J_{\mathrm{eff}}(t)\) be independently measured retained productive uptake over the same accounting window. When the quantities are commensurate and \(J_{\mathrm{raw}}>0\), observed uptake efficiency is
\begin{equation}
    \rho_{\mathrm{obs}}(t)
    =
    \frac{J_{\mathrm{eff}}(t)}{J_{\mathrm{raw}}(t)}.
    \label{eq:effective-exchange}
\end{equation}
In systems with endogenous amplification or noncommensurate input and output units, productive uptake must instead be defined through an explicit conversion or target-response measure.

Capture readiness is measured before or independently of the uptake outcome. CVT may then test a predictor
\begin{equation}
    \widehat{\rho}_{\theta}(t)
    =
    f_{\theta}\!\left(
        G_B(t),G_Q(t),G_C(t),G_S(t),\mathbf w(t)
    \right),
    \label{eq:uptake-predictor}
\end{equation}
where \(\mathbf w(t)\) contains declared exogenous covariates. The predictor may be multiplicative, minimum-gated, nonlinear, or domain-native. Defining \(\rho\) as a product of the same gates and then treating that product as an explanation of uptake is not permitted.

\subsection{Damage Accumulation}

Reserve the symbol \(D(t)\in[0,1]\) for accumulated damage. When raw and effective exchange are commensurate, define unretained delivered load as
\begin{equation}
    L_{\mathrm{u}}(t)
    =
    [J_{\mathrm{raw}}(t)-J_{\mathrm{eff}}(t)]_+.
    \label{eq:unretained-load}
\end{equation}
Otherwise \(L_{\mathrm{u}}\) must be replaced by a directly measured rejection, spillover, waste, or harmful-load quantity.

A candidate damage model is
\begin{equation}
    \dot D
    =
    \gamma L_{\mathrm{u}}
    +
    \delta\sum_i\Phi_i
    +
    \omega R_{\mathrm{risk}}
    -
    \mu(H,S)D,
    \label{eq:damage}
\end{equation}
with projected dynamics preserving \(D\in[0,1]\). The domain-specific risk term \(R_{\mathrm{risk}}\) must be specified before fitting rather than used as a residual container. The repair rate \(\mu(H,S)\) may depend on both learned recovery capacity and currently available reserve. Equation~\eqref{eq:damage} is a candidate mechanism, not a general conservation law.

\subsection{Hosted Transformation as an Independent Outcome}

\begin{definition}[Hosted transformation]
Fix a host state \(x(t)\), a transformation variable \(X(t)\), a domain-native viable set \(\mathcal K_{\mathrm{declared}}\), a transition threshold \(X_{\mathrm{crit}}\), and a follow-up horizon \(\tau_f\geq0\). A transformation is hosted on \([t_0,t_1]\) when
\begin{equation}
    X(t_1)-X(t_0)
    \geq
    X_{\mathrm{crit}},
    \label{eq:transformation-success}
\end{equation}
and
\begin{equation}
    x(t)\in\mathcal K_{\mathrm{declared}}
    \qquad
    \text{for all }t\in[t_0,t_1+\tau_f].
    \label{eq:delayed-host-viability}
\end{equation}
An independently specified damage limit may be added when damage is not already included in \(\mathcal K_{\mathrm{declared}}\).
\end{definition}

The follow-up horizon permits delayed failure to count against apparent success. Gate scores and aggregate CVT scores are predictors or explanatory variables for this outcome; they do not define the outcome they are supposed to predict.

\subsection{Host-Paced Transition as a Model Class}

Let \(u_X(t)\) denote the control or pacing input applied to a transformation. A candidate host-paced policy is
\begin{equation}
    u_X(t)
    =
    u_{\max}
    h_{\theta}\!\left(
        G_B,G_Q,G_C,G_S,D
    \right),
    \qquad
    0\leq h_{\theta}\leq1,
    \label{eq:host-capacity-policy}
\end{equation}
combined with domain-specific transition dynamics
\begin{equation}
    \dot X
    =
    F\!\left(X,u_X,\mathbf w\right).
    \label{eq:host-gated-transition}
\end{equation}
The pacing function \(h_{\theta}\) must not double-count variables already embedded in one another. It should be compared with fixed-rate, capacity-maximizing, and domain-native control policies using the same outcome criteria. CVT does not imply the dimensionally unsupported universal inequality \(\dot X\leq R_{\mathrm{host}}\); any rate bound must be derived in the units and dynamics of the specific application.

\subsection{Minimum Viable Host as a Domain-Specific Threshold}

For a fixed model, forcing severity \(\Lambda\), observation horizon \(\tau\), and control policy \(\pi\), one may estimate a reserve threshold
\begin{equation}
    S(t_0)
    \geq
    S_{\mathrm{crit}}(\Lambda,\tau;\pi,\theta)
    \label{eq:minimum-host-threshold}
\end{equation}
associated with a chosen probability or level of hosted success. This is a domain- and model-specific threshold, not a general theorem. Its form may be nonlinear, history-dependent, or absent. A case that hosts the transformation with no measurable restorative reserve would constrain the claim that reserve is essential in that domain.

\subsection{Local Exchange--Uptake Condition}

Suppose \(J_{\mathrm{raw}}>0\), \(\rho_{\mathrm{obs}}>0\), and
\[
    J_{\mathrm{eff}}
    =
    \rho_{\mathrm{obs}}(J_{\mathrm{raw}})J_{\mathrm{raw}}
\]
locally with differentiable \(\rho_{\mathrm{obs}}\). Then
\begin{equation}
    \frac{dJ_{\mathrm{eff}}}{dJ_{\mathrm{raw}}}
    =
    \rho_{\mathrm{obs}}
    +
    J_{\mathrm{raw}}
    \frac{d\rho_{\mathrm{obs}}}{dJ_{\mathrm{raw}}}.
    \label{eq:exchange-uptake-derivative}
\end{equation}
Effective uptake decreases with additional delivered exchange exactly when
\begin{equation}
    \frac{d\log\rho_{\mathrm{obs}}}{d\log J_{\mathrm{raw}}}
    <
    -1.
    \label{eq:uptake-elasticity-condition}
\end{equation}
This is a calculus condition, not evidence that such degradation occurs. The empirical CVT claim is that some scoped systems enter regimes satisfying Equation~\eqref{eq:uptake-elasticity-condition} because capture readiness, restorative reserve, or another essential mechanism deteriorates under load.

The formal research burden is therefore explicit: define viability independently, publish the raw measurements and normalization maps, estimate uncertainty, compare aggregation and pacing models out of sample, and search deliberately for transformations that remain hosted despite failure of a proposed essential condition.

% CVT Foundations v0.2 — proposed Membrane Illustration replacement
% Review artifact. This section is intended to replace only the current
% membrane material and to stop immediately before the fusion section.

\section{Membrane Illustration: Regulated Transport Across an Interface}
\label{sec:membrane}

A regulated transport interface provides a bounded physical illustration of CVT. It is not presented as the universal mathematical form of the theory. The purpose of the illustration is narrower: to show how delivered flux, interface integrity, pre-existing capture readiness, productive uptake, restorative reserve, and delayed damage can be represented as distinct quantities in a system whose boundary is physically explicit.

\subsection{Established Interface-Transport Form}

Let domains \(A\) and \(B\) meet at an interface \(\Gamma\), and let \(c_A(x,t)\) and \(c_B(x,t)\) denote concentrations or other transport potentials on the two sides. Classical diffusion and membrane-transport models use boundary or interface laws that relate flux to a potential difference and an interface-transfer coefficient \cite{crank1975,kedem1958}. A linear permeability law may be written
\begin{equation}
    J_{\mathrm{raw}}(t)
    =
    \Pi_{\Gamma}(t)
    \bigl(c_B(t)-c_A(t)\bigr)
    \qquad \text{on }\Gamma,
    \label{eq:membrane-raw-flux}
\end{equation}
with the corresponding diffusive boundary flux on side \(A\), under a stated sign convention,
\begin{equation}
    -D_A\nabla c_A\cdot n_A
    =
    J_{\mathrm{raw}}.
    \label{eq:membrane-robin-law}
\end{equation}
Here \(D_A\) is a diffusion coefficient, \(n_A\) is the outward normal, and \(\Pi_{\Gamma}\) is an interface conductance or permeability with units chosen consistently with the transported quantity. Equations~\eqref{eq:membrane-raw-flux}--\eqref{eq:membrane-robin-law} are established transport forms, not CVT inventions. State dependence of \(\Pi_{\Gamma}\) is a modeling choice that must be justified by a physical mechanism or controller.

The interface law identifies what can cross. It does not by itself establish whether the receiving state can capture, retain, or survive the delivered input.

\subsection{Separating Permeability, Flux, Capture, and Uptake}

The membrane illustration distinguishes five quantities that should not be collapsed:
\begin{equation}
    \Pi_{\Gamma}
    \neq
    J_{\mathrm{raw}}
    \neq
    C
    \neq
    \rho_{\mathrm{obs}}
    \neq
    J_{\mathrm{eff}}.
    \label{eq:membrane-separation}
\end{equation}
The permeability \(\Pi_{\Gamma}\) is an interface property or control setting. The raw flux \(J_{\mathrm{raw}}\) is delivered transport. Capture readiness \(C\) is a pre-existing property of the receiving state, estimated before or independently of the uptake outcome. The quantity \(J_{\mathrm{eff}}\) is retained productive uptake measured over a declared accounting window. When raw and effective quantities are commensurate and delivered flux is positive, observed uptake efficiency is
\begin{equation}
    \rho_{\mathrm{obs}}
    =
    \frac{J_{\mathrm{eff}}}{J_{\mathrm{raw}}}.
    \label{eq:membrane-observed-uptake}
\end{equation}
This ratio is an outcome. It must not be used to define the same capture-readiness variable that is then claimed to explain it. In systems with endogenous amplification, storage delay, or noncommensurate input and output units, productive uptake requires an explicit conversion or target-response measure instead of Equation~\eqref{eq:membrane-observed-uptake}.

\subsection{Operational Mapping to \(B,Q,C,S\)}

For a declared receiving host, the membrane illustration maps the four candidate CVT conditions as follows:
\begin{enumerate}[label=(\roman*)]
    \item \textbf{bounded coupling \(B\)}: delivered flux remains within a two-sided admissible band relative to the host's changing handling capacity;
    \item \textbf{preserved identity-defining organization \(Q\)}: the interface and the host organization required by the success criterion remain functionally intact;
    \item \textbf{capture readiness \(C\)}: before the uptake outcome is observed, the receiving state lies in a region from which the specified input can be retained within an admissible basin over the declared horizon;
    \item \textbf{available restorative reserve \(S\)}: measurable recovery headroom remains after transport-induced strain.
\end{enumerate}

Let \(z_B,z_Q,z_C,z_S\) denote domain-native measurements and let \(G_i=g_i(z_i;\vartheta_i)\) be documented normalizations. The candidate membrane viability region is
\begin{equation}
    \mathcal K_{\mathrm{mem}}(t)
    =
    \mathcal K_B(t)
    \cap
    \mathcal K_Q(t)
    \cap
    \mathcal K_C(t)
    \cap
    \mathcal K_S(t).
    \label{eq:membrane-viability-region}
\end{equation}
This intersection is the representation-independent claim. When a scalar score is useful, candidate forms include
\begin{equation}
    \V_{\mathrm{prod,mem}}
    =
    G_BG_QG_CG_S,
    \qquad
    \V_{\min,\mathrm{mem}}
    =
    \min\{G_B,G_Q,G_C,G_S\},
    \label{eq:membrane-candidate-scores}
\end{equation}
and the wider aggregation family in Equation~\eqref{eq:aggregation-family}. No one form is privileged without comparative evidence against an independently declared viability outcome.

\subsection{Transport Regimes and Failure Signatures}

The same interface can enter several distinct regimes.

\paragraph{Insufficient delivery.}
Raw flux remains below the lower edge of the admissible operating band. The host may preserve organization and reserve while failing to receive enough input for the declared transformation. This is a bounded-coupling failure at the low-input side, not evidence that openness is always beneficial.

\paragraph{Overexposure.}
Raw flux exceeds the receiving system's handling band. Damage, rejection, spillover, saturation, or escape from the viable set may increase even when permeability and gradient remain physically large. Overexposure is first a failure of bounded coupling; observed uptake loss may follow but should be measured rather than assumed.

\paragraph{Capture failure.}
The interface delivers input while the receiving state lacks a viable capture path. A low value of \(\rho_{\mathrm{obs}}\) may be consistent with this mechanism, but it does not uniquely identify it: conversion losses, measurement error, transport delay, or an omitted state variable remain alternatives.

\paragraph{Loss of identity-defining organization.}
The interface or host organization required by the declared success criterion fails. Complete merger is a CVT failure only when continued distinction is part of that criterion. If merger is the intended transformation, the host and level of analysis must be redefined rather than treating all loss of boundary as pathological.

\paragraph{Reserve exhaustion.}
Transport remains possible and uptake may remain temporarily positive, but restorative reserve falls below the level required for repair or repeated-load survival. The failure may therefore appear only during a follow-up horizon after the target state has been reached.

These regimes show why permeability alone cannot classify membrane health. The same value of \(\Pi_{\Gamma}\) may be viable or destructive depending on the gradient, receiving state, damage, reserve, observation window, and declared success criterion.

\subsection{Causal State-Dependent Interface Control}

A regulated interface may adapt permeability to estimates of receiving-state condition. A causal policy class is
\begin{equation}
    \Pi_{\Gamma}(t)
    =
    \Pi_{\min}
    +
    \bigl(\Pi_{\max}-\Pi_{\min}\bigr)
    h_{\theta}\!\left(
        \widehat C(t-\Delta t),
        G_S(t-\Delta t),
        D(t-\Delta t),
        z_{\Gamma}(t-\Delta t)
    \right),
    \label{eq:membrane-causal-control}
\end{equation}
where \(0\leq h_{\theta}\leq1\), \(\Delta t\geq0\) is the sensing and actuation delay, and \(z_{\Gamma}\) contains declared interface measurements. A lagged value of \(\rho_{\mathrm{obs}}\) may enter \(z_{\Gamma}\) as feedback from a completed accounting window. Contemporaneous observed uptake cannot be treated as a pre-existing causal input without specifying how it is measured before the control action it is supposed to determine.

Equation~\eqref{eq:membrane-causal-control} is a controller class, not a universal membrane law. It should be compared with fixed permeability, capacity-maximizing control, and domain-native controllers using the same hosted-transformation outcome and follow-up horizon.

\subsection{Minimal Membrane Test}

A reproducible test should declare the host, transported quantity, interface law, viable set, transformation criterion, and follow-up horizon before fitting CVT variables. It should measure pre-input state and reserve, perturb permeability or driving gradient over a stated range, record delivered flux, retained productive uptake, rejection or damage, and recovery, and then compare product, minimum, soft-minimum, and direct-constraint representations out of sample. Deliberate counterexamples should include intended merger, successful transport with negligible measured reserve, high flux without uptake loss, and low capture estimates followed by stable productive retention.

The membrane illustration supports CVT only if the \(B,Q,C,S\) decomposition improves prediction or explanation beyond simpler transport, damage, or state-constraint models. A standard transport model that performs equally well with fewer independently measured quantities would constrain the added value of CVT in this domain.

\subsection{From the Interface Illustration to Fusion}

The next section uses artificial fusion as an illustrative boundary-mediated host problem. The transfer is structural: a plasma--machine system contains interfaces across which heat, particles, waves, fuel, ash, and control influence are transported or regulated while machine viability matters. This does not establish that fusion is literally reducible to Equation~\eqref{eq:membrane-robin-law}, that a tokamak is a biological membrane, or that the membrane illustration captures reactor physics. Those stronger claims require domain-native equations, calibrated parameters, and comparison with established fusion models.

\section{Fusion Instantiation: Hosted Burn}
\label{sec:fusion}

Artificial fusion provides a physically vivid instantiation of CVT. A burning plasma is not merely an energetic state inside a machine. It is a coupled plasma--machine system in which fuel, heat, ash, alpha particles, waves, control signals, and material stresses must be exchanged across an interface. The reactor fails if this exchange is absent, excessive, unreceptive, or non-restorative. It succeeds only if the burn is hosted within a viable exchange regime.

In this section, we interpret the tokamak-like fusion problem as a CVT membrane-host problem. The plasma-machine interface functions as a placental exchange organ: it must nourish the burn without allowing the burn to consume its container.

\subsection{Ignition Access Is Not Hosted Burn}

Fusion research often treats ignition as the decisive transition. In CVT terms, ignition is only the beginning of the host problem. Before ignition, the machine primarily drives the plasma. After alpha heating becomes significant, the plasma becomes an endogenous source of heat and load.

Let
\begin{equation}
    \chialpha
    =
    \frac{P_{\alpha}}{P_{\alpha}+P_{\mathrm{aux}}},
\end{equation}
where \(P_{\alpha}\) is alpha-particle self-heating and \(P_{\mathrm{aux}}\) is externally supplied heating. When \(\chialpha\approx 0\), the plasma is externally driven. When \(\chialpha\to 1\), heating is alpha-dominated.

Ignition access concerns the transition
\[
    \chialpha \uparrow.
\]
Hosted burn concerns whether this transition can occur while the host remains viable:
\begin{equation}
    \V(t)>\V_{\mathrm{crit}},\qquad
    \rho(t)>\rho_{\mathrm{crit}},\qquad
    D(t)<D_{\mathrm{crit}}.
\end{equation}

Thus CVT distinguishes:
\[
    \text{ignition access}
    \neq
    \text{hosted burn}.
\]

A reactor may reach alpha-dominant conditions while losing receptivity and saturating damage. In that case, ignition has occurred in a narrow sense, but hosted burn has failed.

\subsection{The Placental Exchange Interface}

We define the placental exchange interface as the coupled system of machine functions that allow the burn to persist without destructive contact. These include:
\begin{itemize}
    \item fuel delivery,
    \item helium ash removal,
    \item heat exhaust,
    \item blanket and shielding,
    \item coolant throughput,
    \item material tolerance,
    \item magnetic confinement,
    \item spectral and wave control,
    \item diagnostics and feedback authority.
\end{itemize}

This interface is placental in the structural sense: it mediates exchange between a developing energetic process and the host that carries it. The term does not import biological detail into fusion physics. The claim is structural. A burn requires exchange. The host requires separation. The interface must allow one without destroying the other.

Let the relevant gates be
\[
G_{\mathrm{burn}},\quad
G_{\mathrm{machine}},\quad
G_{\mathrm{exchange}},\quad
G_{\mathrm{separation}},\quad
G_{\mathrm{restoration}},\quad
G_{\mathrm{spectral}}.
\]

Total placental viability is
\begin{equation}
    \V_P
    =
    G_{\mathrm{burn}}
    G_{\mathrm{machine}}
    G_{\mathrm{exchange}}
    G_{\mathrm{separation}}
    G_{\mathrm{restoration}}
    G_{\mathrm{spectral}}.
    \label{eq:placental-viability}
\end{equation}

The product form states that no single subsystem can compensate indefinitely for collapse of another. A strong magnet does not compensate for failed heat exhaust. Strong fueling does not compensate for ash accumulation. Strong burn does not compensate for loss of machine integrity.

\subsection{Receptivity and Effective Exchange}

In the fusion instantiation, raw exchange is not sufficient. The machine may deliver fuel, heat, current drive, or wave power, but these inputs are only useful if the coupled plasma-machine state can receive them.

We write:
\begin{equation}
    \Jraw = \Pi \Delta,
\end{equation}
where \(\Pi\) is exchange capacity and \(\Delta\) is the relevant driving gradient. Effective exchange is
\begin{equation}
    \Jeff = \rho \Pi \Delta.
\end{equation}

A minimal fusion receptivity functional is
\begin{equation}
    \rho
    =
    G_{\mathrm{burn}}
    G_{\mathrm{machine}}
    G_{\mathrm{separation}}
    G_{\mathrm{restoration}}
    G_{\mathrm{spectral}}
    (1-D).
\end{equation}

The exchange gate itself is excluded from this expression because receptivity asks whether the host can receive exchange, not whether the exchange channel exists.

This yields a central CVT claim:

\begin{quote}
A reactor may possess exchange capacity while lacking exchange receptivity.
\end{quote}

When \(\rho\to 0\), effective exchange vanishes even if raw exchange remains high:
\[
    \Jeff \to 0
    \qquad
    \text{while}
    \qquad
    \Jraw>0.
\]

In this case, delivered exchange becomes load.

\subsection{Post-Ignition Throttle Transition}

The transition from externally driven plasma to alpha-dominant burn is analogous to the transition from takeoff thrust to climb thrust in flight. Takeoff thrust is useful for leaving the runway, but it is not the sustainable flight regime. Similarly, ignition-access control may be useful before alpha heating begins, but destructive after the plasma becomes an endogenous heat source.

A post-ignition throttle transition reduces or reshapes external drive as \(\chialpha\) rises. A simple form is
\begin{equation}
    P_{\mathrm{aux}}(t)
    =
    P_{\mathrm{takeoff}}(1-\chialpha)
    +
    P_{\mathrm{climb}}\rho\chialpha.
\end{equation}

However, auxiliary throttle-back alone is not sufficient if alpha dominance continues to rise on a fixed schedule. CVT requires that the handoff itself be gated by host capacity:
\begin{equation}
    \dot{\chialpha}
    =
    r_{\alpha}R_{\mathrm{host}}(1-\chialpha)
    -
    s_{\alpha}(1-R_{\mathrm{host}})\chialpha.
    \label{eq:alpha-pacing}
\end{equation}

A minimal host capacity functional is
\begin{equation}
    R_{\mathrm{host}}
    =
    \rho(1-D)
    G_{\mathrm{spectral}}
    G_{\mathrm{restoration}}.
\end{equation}

Equation~\eqref{eq:alpha-pacing} states that alpha dominance rises when the host is receptive and slows or retreats when host viability deteriorates. This is the formal version of post-ignition pacing:

\[
    \frac{d\chialpha}{dt}
    \leq
    R_{\mathrm{host}}(t).
\]

\subsection{Simulation Result 1: Control Philosophy}

A structural toy model compares three controllers: capacity-maximizing, receptivity-gated, and spectral-aware. The model is not a tokamak physics code. It does not solve magnetohydrodynamics, gyrokinetics, alpha-particle transport, impurity dynamics, or scrape-off-layer physics. Its purpose is to test the internal architecture of CVT.

The capacity-maximizing controller delivers the largest raw exchange but collapses receptivity. In one calibration:
\[
    \overline{J}_{\mathrm{raw}}=0.647,
    \qquad
    \overline{J}_{\mathrm{eff}}=0.002,
    \qquad
    D(T)=1,
    \qquad
    \rho(T)=0.
\]

The spectral-aware controller uses less raw exchange but preserves receptivity:
\[
    \overline{J}_{\mathrm{raw}}=0.049,
    \qquad
    \overline{J}_{\mathrm{eff}}=0.010,
    \qquad
    D(T)=0.104,
    \qquad
    \rho(T)=0.259.
\]

Thus the controller that pushes hardest performs worst. The controller that preserves receptivity receives more useful exchange.

\begin{quote}
More delivered exchange is not the same as more received exchange.
\end{quote}

\subsection{Simulation Result 2: Ablation}

The ablation study separates the model's mechanisms. Removing receptivity gating reproduces the overdrive failure:
\[
    \overline{J}_{\mathrm{raw}}=0.538,
    \qquad
    \overline{J}_{\mathrm{eff}}=0.002,
    \qquad
    D(T)=1,
    \qquad
    \rho(T)=0.
\]

Removing maturity avoids immediate collapse but prevents strong recovery:
\[
    \V_P(T)=0.046,
    \qquad
    P(T)=0.
\]

Removing coupled degradation slightly improves performance, as expected, because violations no longer propagate across margins. This confirms that the coupling term represents a real cascade cost.

The ablation supports the following mechanism map:
\begin{center}
\begin{tabular}{ll}
\toprule
Mechanism & Function \\
\midrule
Product gates & Nonlinear fragility \\
Coupled degradation & Cascade and correlated margin loss \\
Receptivity gating & Prevention of overdrive damage \\
Maturity & Recovery after shock \\
Spectral awareness & Protection under multimode stress \\
\bottomrule
\end{tabular}
\end{center}

\subsection{Simulation Result 3: Ignition Survival}

When alpha dominance is forced toward unity on a fixed schedule, all controllers reach
\[
    \chialpha(T)=1.
\]

However, the host fails:
\[
    D(T)=1,
    \qquad
    \rho(T)=0.
\]

Thus the system reaches alpha dominance but loses the capacity to receive it. This motivates the distinction between ignition and hosted burn. In CVT terms, this is over-ignition: the transition succeeds energetically but fails as hosted transformation.

\subsection{Simulation Result 4: Throttle Transition}

The throttle-transition experiment compares fixed alpha ramp, auxiliary throttle-back, and host-gated alpha pacing.

Auxiliary throttle-back reduces external drive:
\[
    \overline{P}_{\mathrm{aux}}=0.260
\]
compared with
\[
    \overline{P}_{\mathrm{aux}}=0.836
\]
in the fixed-ramp case. Nevertheless, damage still saturates and receptivity collapses when alpha dominance rises on a fixed schedule.

Host-gated alpha pacing behaves differently:
\[
    \V_P(T)=0.186,
    \qquad
    P(T)=0.533,
    \qquad
    D(T)=0.283,
    \qquad
    \rho(T)=0.148.
\]

The host remains alive, but alpha dominance is limited:
\[
    \chialpha(T)=0.144.
\]

This shows that pacing preserves the host but may prevent full alpha dominance unless sufficient reserve exists.

\begin{quote}
Auxiliary throttle-back reduces the push; host-gated alpha pacing regulates the handoff.
\end{quote}

\subsection{Simulation Result 5: Minimum Viable Host}

To identify the reserve required for hosted burn, we introduce placental reserve \(R_P\), representing combined heat exhaust, material tolerance, coolant throughput, ash removal, spectral damping, diagnostic fidelity, control authority, magnetic margin, and thermal inertia.

In the model, reserve reduces effective alpha load and increases host transition capacity. We sweep \(R_P\) against alpha-load severity \(\Lambda_{\alpha}\).

Hosted burn is defined by the simultaneous conditions:
\begin{equation}
    \chialpha>0.75,
    \qquad
    D<0.60,
    \qquad
    \rho>0.05,
    \qquad
    \V_P>0.05.
\end{equation}

The sweep identifies a hosted-burn boundary. In one calibration, the strict threshold is approximately:
\begin{equation}
    R_{P,\mathrm{crit}}
    \approx
    6.25+5.5\Lambda_{\alpha}.
    \label{eq:reserve-curve}
\end{equation}

At low alpha-load severity,
\[
    \Lambda_{\alpha}=0.05,
\]
hosted burn begins near
\[
    R_P=6.5.
\]

At higher alpha-load severity,
\[
    \Lambda_{\alpha}=0.50,
\]
the threshold rises to approximately
\[
    R_P=9.0.
\]

A high-reserve hosted-burn case gives:
\[
    R_P=14,
    \qquad
    \Lambda_{\alpha}=0.05,
\]
with
\[
    \chialpha=0.998,
    \qquad
    D=0.085,
    \qquad
    \rho=0.306,
    \qquad
    \V_P=0.321,
    \qquad
    P=0.863.
\]

This is not merely non-collapse. It is a mature hosted-burn state in the toy model: alpha dominance is nearly complete, damage is low, receptivity remains positive, and placental maturity is high.

\subsection{Minimum Viable Host Theorem}

\begin{theorem}[Minimum Viable Host, toy-model form]
For a host-gated alpha pacing system with alpha-load severity \(\Lambda_{\alpha}\), there exists a critical placental reserve threshold
\[
    R_{P,\mathrm{crit}}(\Lambda_{\alpha})
\]
such that hosted burn occurs only when
\[
    R_P\geq R_{P,\mathrm{crit}}(\Lambda_{\alpha}).
\]
Hosted burn requires simultaneous satisfaction of
\[
    \chialpha>\chi_{\mathrm{host}},
    \qquad
    D<D_{\mathrm{crit}},
    \qquad
    \rho>\rho_{\mathrm{crit}},
    \qquad
    \V_P>\V_{\mathrm{crit}}.
\]
\end{theorem}

\begin{remark}
In the calibration reported above,
\[
    \chi_{\mathrm{host}}=0.75,
    \qquad
    D_{\mathrm{crit}}=0.60,
    \qquad
    \rho_{\mathrm{crit}}=0.05,
    \qquad
    \V_{\mathrm{crit}}=0.05,
\]
and the observed threshold is approximately linear:
\[
    R_{P,\mathrm{crit}}
    \approx
    6.25+5.5\Lambda_{\alpha}.
\]
This should not be interpreted as a calibrated engineering law. It is a structural result from a toy model: hosted burn requires reserve sufficient to keep the exchange interface receptive under alpha-dominant load.
\end{remark}

\subsection{Interpretation}

The fusion instantiation supports the broader CVT claim. A system may be capable of initiating a transformation without being capable of hosting it. Ignition access is not equivalent to hosted burn. Auxiliary throttle-back is not equivalent to handoff pacing. Raw exchange is not equivalent to effective exchange.

The reactor becomes an artificial star not when alpha heating begins, but when alpha heating can be handed off to a host whose exchange interface remains viable.

In ordinary language:

\begin{quote}
Ignition is lift-off; hosted burn is flight; placental reserve is the runway plus the wings.
\end{quote}

\begin{thebibliography}{9}

\bibitem{aubin2011}
J.-P. Aubin, A. M. Bayen, and P. Saint-Pierre,
\emph{Viability Theory: New Directions}, 2nd ed.
Springer, 2011. doi:10.1007/978-3-642-16684-6.

\bibitem{maturana1980}
H. R. Maturana and F. J. Varela,
\emph{Autopoiesis and Cognition: The Realization of the Living}.
D. Reidel, 1980. doi:10.1007/978-94-009-8947-4.

\bibitem{vanderschaft2014}
A. van der Schaft and D. Jeltsema,
\emph{Port-Hamiltonian Systems Theory: An Introductory Overview}.
Now Publishers, 2014. doi:10.1561/2600000002.

\bibitem{holling1973}
C. S. Holling,
``Resilience and Stability of Ecological Systems,''
\emph{Annual Review of Ecology and Systematics}, vol. 4, pp. 1--23, 1973.
doi:10.1146/annurev.es.04.110173.000245.

\bibitem{menck2013}
P. J. Menck, J. Heitzig, N. Marwan, and J. Kurths,
``How basin stability complements the linear-stability paradigm,''
\emph{Nature Physics}, vol. 9, pp. 89--92, 2013.
doi:10.1038/nphys2516.

\bibitem{crank1975}
J. Crank, \emph{The Mathematics of Diffusion}, 2nd ed.
Clarendon Press, Oxford, 1975. ISBN: 978-0-19-853344-3.

\bibitem{kedem1958}
O. Kedem and A. Katchalsky,
``Thermodynamic analysis of the permeability of biological membranes to non-electrolytes,''
\emph{Biochimica et Biophysica Acta}, vol. 27, no. 2, pp. 229--246, 1958.
doi:10.1016/0006-3002(58)90330-5.

\end{thebibliography}

\end{document}

